\chapter{Cross-Cloud Service Equivalence}
\index{cross-cloud equivalence}\index{service mapping}\index{multi-cloud}%
This appendix is a consolidated reference that maps conceptually equivalent services and
ideas across the seven platforms studied in this book series: Amazon Web Services
(AWS), Microsoft Azure, Microsoft Fabric, Google Cloud Platform (GCP), MongoDB,
Snowflake, and Databricks. The names differ, but the underlying data-engineering concepts---durable
object storage, tiered lifecycle economics, immutability, managed relational engines,
elastic compute, and disciplined data organization---recur everywhere.

\begin{notebox}
This appendix grew with each chapter. Every chapter contributed the equivalence
tables introduced in its cross-cloud callout, so the book now holds
a single consolidated Rosetta Stone for moving fluently between clouds. Where a clean
one-to-one mapping does not exist, the prose notes the closest analog and the caveat.
Databricks, which runs on all three hyperscalers, receives both per-domain treatment and a
consolidated at-a-glance table at the end.
\end{notebox}

\section{Object and Data-Lake Storage}
\index{object storage!equivalence}\index{storage classes!equivalence}%
The foundational layer covered in Chapter~1. Every platform offers durable, HTTP-addressable
object storage with tiered pricing, lifecycle transitions, versioning, and write-once-read-many
(WORM) controls.

\begin{center}
\footnotesize
\begin{tabular}{@{}p{2.8cm}p{3.2cm}p{3.2cm}p{3.2cm}@{}}
\toprule
\textbf{Capability} & \textbf{AWS} & \textbf{Azure} & \textbf{GCP} \\
\midrule
Object storage & Amazon S3 & ADLS Gen2 / Blob & Cloud Storage \\
Hot tier & S3 Standard & Hot & Standard \\
Infrequent access & S3 Standard-IA & Cool & Nearline \\
Archive (instant) & Glacier Instant Retrieval & Cold & Coldline \\
Deep archive & Glacier Deep Archive & Archive & Archive \\
Lifecycle rules & Lifecycle policies & Lifecycle management & Object Lifecycle Mgmt \\
Versioning & S3 Versioning & Blob versioning & Object Versioning \\
Immutability (WORM) & Object Lock & Immutable blob storage & Bucket Lock \\
Cross-region copy & CRR / SRR & Object replication & Dual/multi-region + Transfer \\
Encryption at rest & SSE-S3 / SSE-KMS & SSE + CMK (Key Vault) & Google-managed / CMEK \\
\bottomrule
\end{tabular}
\end{center}

\noindent The three remaining platforms approach storage differently.
\textbf{Microsoft Fabric} exposes a single logical lake, \emph{OneLake}, and uses
\emph{shortcuts} to virtualize S3, ADLS Gen2, and Google Cloud Storage in place without copying.
\textbf{Snowflake} abstracts storage entirely: table data lives in cloud object storage as
immutable, columnar \emph{micro-partitions} that the service manages on your behalf.
\textbf{MongoDB} persists documents through the \emph{WiredTiger} storage engine, balancing an
internal cache against the operating-system page cache rather than exposing storage tiers.

\section{Managed Relational Databases}
\index{relational databases!equivalence}%
The focus of Chapter~2 for the three hyperscalers. Each provides a fully managed relational
engine, a cloud-native high-performance variant, high availability, read replicas, and a
migration service.

\begin{center}
\footnotesize
\begin{tabular}{@{}p{2.8cm}p{3.2cm}p{3.2cm}p{3.2cm}@{}}
\toprule
\textbf{Capability} & \textbf{AWS} & \textbf{Azure} & \textbf{GCP} \\
\midrule
Managed relational & Amazon RDS & Azure SQL Database & Cloud SQL \\
Cloud-native engine & Amazon Aurora & SQL Hyperscale & AlloyDB \\
Managed PostgreSQL & RDS / Aurora PostgreSQL & Azure DB for PostgreSQL & Cloud SQL / AlloyDB \\
High availability & Multi-AZ & Zone-redundant HA & Regional HA \\
Read scaling & Read Replicas & Read replicas / geo-replicas & Read Replicas \\
Global / geo & Aurora Global Database & Active geo-replication & Cross-region replicas \\
Serverless compute & Aurora Serverless v2 & Azure SQL Serverless & AlloyDB (elastic) \\
Migration service & DMS + SCT & Azure Database Migration Svc & Database Migration Svc \\
\bottomrule
\end{tabular}
\end{center}

\noindent \textbf{Snowflake}, \textbf{Fabric}, and \textbf{MongoDB} are not managed relational
engines, but they solve overlapping problems: Snowflake and Fabric provide SQL analytics over
warehouse/lakehouse storage, while MongoDB provides a distributed document store whose schema
design (Chapter~2) replaces rigid relational normalization with intentional, access-driven modeling.

\section{Elastic Compute and Analytical Warehouses}
\index{elastic compute!equivalence}\index{virtual warehouse!equivalence}%
Snowflake's Chapter~2 topic---separating compute from storage and scaling it elastically---has a
counterpart in every analytics platform.

\begin{center}
\footnotesize
\begin{tabular}{@{}p{2.8cm}p{3.2cm}p{3.2cm}p{3.2cm}@{}}
\toprule
\textbf{Concept} & \textbf{Snowflake} & \textbf{AWS} & \textbf{Azure / GCP} \\
\midrule
Compute unit & Virtual Warehouse & Redshift RA3 / Serverless & Synapse pool / BigQuery slots \\
Elastic concurrency & Multi-cluster warehouse & Concurrency scaling & Autoscale / on-demand \\
Pause when idle & Auto-suspend / resume & Pause / resume & Auto-pause / serverless \\
Billing unit & Credits & RPU / node-hours & DWU / slot-hours \\
Result reuse & Result cache & Result cache & Result / BI cache \\
\bottomrule
\end{tabular}
\end{center}

\noindent \textbf{Microsoft Fabric} meters all workloads against a shared pool of Capacity Units
(CU) on an F-SKU, an approach closer to a single elastic budget than to per-engine warehouses.
\textbf{MongoDB Atlas} scales through cluster tiers with optional auto-scaling of compute and storage.

\section{Data Organization and the Medallion Pattern}
\index{medallion architecture!equivalence}%
Fabric's Chapter~2 topic---the Bronze/Silver/Gold medallion---is a portable design pattern rather
than a product. Every platform implements the same progressive-refinement idea with its own primitives.

\begin{center}
\footnotesize
\begin{tabular}{@{}p{2.6cm}p{3.4cm}p{3.4cm}p{3.0cm}@{}}
\toprule
\textbf{Layer} & \textbf{Fabric} & \textbf{Data lake (AWS/Azure/GCP)} & \textbf{Snowflake} \\
\midrule
Bronze (raw) & Bronze Lakehouse & \texttt{raw/} prefix (S3/ADLS/GCS) & RAW schema / stage \\
Silver (validated) & Silver Lakehouse & \texttt{validated/} prefix & STAGING schema \\
Gold (curated) & Gold Lakehouse & \texttt{curated/} prefix & ANALYTICS schema \\
Storage format & Delta-Parquet & Parquet / Delta / Iceberg & Micro-partitions \\
\bottomrule
\end{tabular}
\end{center}

\noindent \textbf{MongoDB} realizes the same separation with distinct collections or databases
(for example \texttt{raw\_events}, \texttt{validated}, \texttt{curated}) governed by schema-design
patterns rather than a lake layout.

\section{Document Data Modeling}
\index{document modeling!equivalence}%
MongoDB's Chapter~2 schema-design patterns generalize to every document and wide-column store.

\begin{center}
\footnotesize
\begin{tabular}{@{}p{3.0cm}p{3.2cm}p{3.2cm}p{3.0cm}@{}}
\toprule
\textbf{Concept} & \textbf{MongoDB} & \textbf{Amazon DynamoDB} & \textbf{Azure Cosmos DB} \\
\midrule
Primary access unit & Document / collection & Item / table & Item / container \\
Partitioning & Shard key & Partition key & Partition key \\
One-table modeling & Embedding + patterns & Single-table design & Embedding + partitioning \\
Time-series grouping & Bucket Pattern & Composite sort keys & Bucketing by key \\
Optional attributes & Attribute Pattern & Sparse attributes & Sparse properties \\
\bottomrule
\end{tabular}
\end{center}

\noindent Google Cloud's closest analogs are \textbf{Firestore} and \textbf{Bigtable}. In
\textbf{Snowflake} and \textbf{Fabric}, semi-structured documents are stored in \texttt{VARIANT}
(Snowflake) or JSON columns and queried with path expressions rather than modeled as collections.

\section{Globally Distributed and NoSQL Databases}
\index{NoSQL!equivalence}\index{distributed databases!equivalence}%
Chapter~3 branches by platform into high-scale NoSQL (AWS DynamoDB, Azure Cosmos DB), globally
consistent NewSQL (GCP Cloud Spanner), and advanced document modeling (MongoDB). They all answer
the same question: how to scale a database horizontally across partitions and regions.

\begin{center}
\footnotesize
\begin{tabular}{@{}p{2.8cm}p{2.9cm}p{2.9cm}p{2.9cm}p{2.4cm}@{}}
\toprule
\textbf{Concept} & \textbf{DynamoDB} & \textbf{Cosmos DB} & \textbf{Cloud Spanner} & \textbf{MongoDB} \\
\midrule
Model & Key-value / doc & Multi-model & Distributed SQL & Document \\
Partitioning & Partition key & Partition key & Split by key range & Shard key \\
Throughput unit & RCU / WCU & Request Units (RU) & Nodes / processing units & Cluster tier \\
Secondary indexes & GSI / LSI & Automatic indexing & Secondary indexes & Indexes \\
Change stream & DynamoDB Streams & Change Feed & Change streams & Change streams \\
Expiry & TTL & TTL & (manual) & TTL index \\
Multi-region & Global Tables & Multi-region writes & Multi-region (synchronous) & Zones / global clusters \\
Consistency & Eventual / strong & 5 tunable levels & External (TrueTime) & Tunable read/write concern \\
\bottomrule
\end{tabular}
\end{center}

\noindent \textbf{Snowflake} and \textbf{Fabric} are analytical platforms rather than OLTP key-value
stores; they ingest from these operational databases and store semi-structured payloads in
\texttt{VARIANT}/JSON columns. Google's dedicated NoSQL analogs are \textbf{Firestore} (document)
and \textbf{Bigtable} (wide-column).

\section{Managed Apache Spark}
\index{Apache Spark!equivalence}%
Microsoft Fabric's Chapter~3 topic---managed Spark pools and notebooks---has a first-class
equivalent on every cloud.

\begin{center}
\footnotesize
\begin{tabular}{@{}p{2.8cm}p{3.2cm}p{3.2cm}p{3.2cm}@{}}
\toprule
\textbf{Capability} & \textbf{Fabric} & \textbf{AWS} & \textbf{Azure / GCP} \\
\midrule
Managed Spark & Starter / Custom pools & Glue / EMR / EMR Serverless & Synapse Spark / Databricks; Dataproc \\
Notebook UX & Fabric Notebooks & Glue Studio / EMR Studio & Synapse / Databricks; Vertex AI \\
Serverless option & Autoscaling pools & Glue / EMR Serverless & Synapse serverless; Dataproc Serverless \\
Lake format & Delta on OneLake & Delta / Iceberg / Hudi & Delta; Delta / Iceberg \\
\bottomrule
\end{tabular}
\end{center}

\noindent \textbf{Snowflake} offers \textbf{Snowpark} for DataFrame-style processing, and
\textbf{MongoDB} integrates through the Spark Connector.

\section{Query Acceleration and Materialized Views}
\index{materialized views!equivalence}\index{query acceleration!equivalence}%
Snowflake's Chapter~3 acceleration tools---materialized views, the Search Optimization Service, and
transient/temporary tables---map to comparable features in every warehouse.

\begin{center}
\footnotesize
\begin{tabular}{@{}p{2.8cm}p{3.0cm}p{3.0cm}p{3.4cm}@{}}
\toprule
\textbf{Feature} & \textbf{Snowflake} & \textbf{AWS Redshift} & \textbf{Azure Synapse / GCP BigQuery} \\
\midrule
Materialized views & Materialized Views & Materialized views & Materialized views \\
Point-lookup accel. & Search Optimization Svc & Sort keys / zone maps & (indexing) / search indexes \\
Result reuse & Result cache & Result cache & Result-set cache / BI Engine \\
Ephemeral tables & Transient / temporary & Temp tables & Temp tables \\
\bottomrule
\end{tabular}
\end{center}

\noindent \textbf{Fabric} accelerates reads through Direct Lake and its Warehouse engine, while
\textbf{MongoDB} relies on secondary and compound indexes plus the in-memory working set.

\section{In-Memory Caching and Hybrid Ingestion}
\index{caching!equivalence}\index{in-memory cache!equivalence}%
The AWS and Azure Chapter~4 topic---low-latency caching in front of a durable store, fed by hybrid
(edge-to-cloud) ingestion---has a managed Redis-compatible equivalent on every cloud.

\begin{center}
\footnotesize
\begin{tabular}{@{}p{2.8cm}p{3.0cm}p{3.0cm}p{3.4cm}@{}}
\toprule
\textbf{Capability} & \textbf{AWS} & \textbf{Azure} & \textbf{GCP / others} \\
\midrule
Managed cache & ElastiCache (Redis/Memcached) & Cache for Redis & Memorystore (Redis/Memcached) \\
Real-time API layer & AppSync / API Gateway & API Management / SignalR & API Gateway / Apigee \\
Hybrid ingestion & DataSync / Snow family / IoT Core & Data Box / IoT Hub / Arc & Transfer Appliance / IoT Core \\
Streaming intake & Kinesis / MSK & Event Hubs & Pub/Sub \\
\bottomrule
\end{tabular}
\end{center}

\noindent \textbf{Snowflake} and \textbf{Fabric} rely on result and Direct~Lake caches rather than a
standalone Redis tier, while \textbf{MongoDB} serves hot data from its in-memory working set and
Atlas can pair with an external cache.

\section{Wide-Column and Key-Value NoSQL}
\index{wide-column!equivalence}\index{Bigtable!equivalence}%
GCP's Chapter~4 topic---Cloud Bigtable---is a wide-column, high-throughput, low-latency NoSQL store.
Its peers span the managed and open-source ecosystems.

\begin{center}
\footnotesize
\begin{tabular}{@{}p{2.8cm}p{3.0cm}p{3.0cm}p{3.4cm}@{}}
\toprule
\textbf{Characteristic} & \textbf{GCP Bigtable} & \textbf{AWS} & \textbf{Azure / open source} \\
\midrule
Data model & Wide-column & DynamoDB (key-value/doc) & Cosmos DB (multi-model) \\
HBase-compatible & HBase API & Keyspaces (Cassandra) & Managed Instance for Cassandra \\
Scaling unit & Nodes / clusters & On-demand / provisioned & RU/s (Cosmos) \\
Best fit & Time-series, IoT, ad-tech & Serverless key-value & Global multi-model \\
\bottomrule
\end{tabular}
\end{center}

\noindent \textbf{Snowflake} and \textbf{Fabric} address these workloads analytically rather than as
operational key-value stores; \textbf{MongoDB} covers the document end of the same NoSQL spectrum.

\section{Data Protection: Cloning, Time Travel, and Migration}
\index{cloning!equivalence}\index{time travel!equivalence}\index{point-in-time recovery!equivalence}%
Snowflake's Chapter~4 topic---zero-copy cloning, Time Travel, and Fail-safe---plus MongoDB's
relational-to-document migration map to data-protection and migration tooling across the clouds.

\begin{center}
\footnotesize
\begin{tabular}{@{}p{2.8cm}p{3.0cm}p{3.0cm}p{3.4cm}@{}}
\toprule
\textbf{Capability} & \textbf{Snowflake} & \textbf{AWS / Azure} & \textbf{GCP / Fabric} \\
\midrule
Instant clone & Zero-copy clone & Aurora clone / SQL DB copy & BigQuery table clone; OneLake shortcut \\
Point-in-time recovery & Time Travel & PITR (RDS/SQL DB backups) & PITR; Delta time travel \\
Retention safety net & Fail-safe (7 days) & Automated backups & Backups; Delta VACUUM window \\
Schema migration & --- & DMS / SCT & Database Migration Service \\
\bottomrule
\end{tabular}
\end{center}

\noindent \textbf{MongoDB} migrations use \textbf{Relational Migrator} plus the embed-vs-reference
patterns from Chapter~4, and \textbf{Delta Lake} exposes historical versions through
\texttt{VERSION AS OF} time-travel queries.

\section{Serverless Analytical Warehouses}
\index{warehouse!equivalence}\index{BigQuery!equivalence}%
Chapter~5's topic---BigQuery's serverless, columnar warehouse---is the analytical heart of every
platform. Pricing mechanics and tuning surfaces differ; decoupled storage and compute do not.

\begin{center}
\footnotesize
\begin{tabular}{@{}p{2.8cm}p{3.0cm}p{3.0cm}p{3.4cm}@{}}
\toprule
\textbf{Capability} & \textbf{GCP BigQuery} & \textbf{AWS Redshift} & \textbf{Azure / Databricks} \\
\midrule
Warehouse engine & Dremel over Colossus & RA3 managed storage & Synapse / SQL warehouse (Photon) \\
Unit of compute & Slots & RPUs / node hours & DWU / DBU \\
On-demand pricing & Per TiB scanned & Serverless RPUs & Serverless SQL per TB / serverless DBU \\
Capacity pricing & Editions + reservations & Provisioned clusters & Provisioned pools / pro \& classic warehouses \\
Result acceleration & BI Engine & Materialized views, AQUA & Result-set cache / DBSQL cache \\
Nested data & ARRAY and STRUCT & SUPER (semi-structured) & OPENJSON / Spark nested types \\
\bottomrule
\end{tabular}
\end{center}

\noindent \textbf{Snowflake} is the closest philosophical cousin (virtual warehouses billed per
credit); \textbf{Fabric} meters everything against capacity units. Databricks SQL warehouses query
the \emph{same} Delta tables pipelines write---no copy into a proprietary store.

\section{Query Optimization and Physical Layout}
\index{partitioning!equivalence}\index{clustering!equivalence}%
Chapter~6's two levers---eliminate data before it is read, precompute what is read often---exist
on every analytical platform under different names.

\begin{center}
\footnotesize
\begin{tabular}{@{}p{2.8cm}p{3.0cm}p{3.0cm}p{3.4cm}@{}}
\toprule
\textbf{Capability} & \textbf{GCP BigQuery} & \textbf{Snowflake} & \textbf{Databricks (Delta)} \\
\midrule
Coarse pruning & Partitioning & Micro-partition pruning & Partitioning / file stats \\
Fine-grained ordering & Clustering (4 columns) & Clustering keys & \texttt{ZORDER BY} / liquid clustering \\
Compaction & Managed automatically & Auto reclustering & \texttt{OPTIMIZE} \\
Pre-aggregation & Materialized views & Materialized views & Materialized views \\
In-memory tier & BI Engine & Result / warehouse cache & Disk cache / Photon \\
Auto-maintenance & Autoclustering & Automatic & Predictive optimization \\
\bottomrule
\end{tabular}
\end{center}

\noindent \textbf{AWS Redshift} exposes sort and dist keys plus zone maps; \textbf{Fabric} adds
V-Order; \textbf{MongoDB} answers the same question with indexes and covered queries. The physics
is shared: touch less data per query.

\section{Federated and Cross-Cloud Analytics}
\index{federation!equivalence}\index{cross-cloud analytics!equivalence}%
Chapter~7's questions---``can I query data I have not loaded?'' and ``can I query data in another
cloud?''---have answers everywhere, with different mechanics and egress bills.

\begin{center}
\footnotesize
\begin{tabular}{@{}p{2.8cm}p{3.0cm}p{3.0cm}p{3.4cm}@{}}
\toprule
\textbf{Capability} & \textbf{GCP BigQuery} & \textbf{AWS} & \textbf{Azure / Databricks} \\
\midrule
Query files in place & External tables / BigLake & Redshift Spectrum / Athena & Synapse serverless / external tables, Delta over S3/ADLS/GCS \\
Cross-cloud analytics & BigQuery Omni & Athena federated queries & Fabric shortcuts / multi-cloud workspaces \\
Geospatial & GEOGRAPHY (S2) & GEOMETRY / GEOGRAPHY & KQL geo / H3, spatial SQL \\
Full-text search & SEARCH indexes & OpenSearch integration & AI Search / Lakehouse search \\
\bottomrule
\end{tabular}
\end{center}

\noindent \textbf{Snowflake} queries external and Iceberg tables in customer storage;
\textbf{MongoDB Atlas Data Federation} reaches into S3 and Atlas clusters. The universal lesson:
federation moves compute to data only when the engine supports it---otherwise bytes travel and the
egress bill follows.

\section{Lakehouse Governance, Sharing, and Clean Rooms}
\index{lakehouse!equivalence}\index{data sharing!equivalence}\index{clean rooms!equivalence}%
Chapter~8's lakehouse pattern---open files in object storage with warehouse-grade security on
top---exists on every platform, with different names for the unification layer.

\begin{center}
\footnotesize
\begin{tabular}{@{}p{2.8cm}p{3.0cm}p{3.0cm}p{3.4cm}@{}}
\toprule
\textbf{Capability} & \textbf{GCP} & \textbf{AWS} & \textbf{Databricks / Fabric} \\
\midrule
Lakehouse tables & BigLake (+ Iceberg) & Lake Formation governed tables & Delta + Unity Catalog / OneLake lakehouse \\
Unstructured indexing & Object Tables & Athena + external metadata & Volumes + AI/BI / OneLake files \\
Row-level security & Row access policies & Lake Formation row filters & UC row filters / SQL RLS \\
Column-level security & Policy tags & Lake Formation column grants & UC column masks / Purview \\
Data sharing & Analytics Hub & Data Exchange / Redshift sharing & Delta Sharing / external sharing \\
Data clean rooms & Analytics Hub clean rooms & AWS Clean Rooms & UC clean rooms / confidential \\
\bottomrule
\end{tabular}
\end{center}

\noindent \textbf{Snowflake} delivers the same story through external tables, masking policies, and
Marketplace listings; \textbf{MongoDB}'s analog is field-level encryption and role-based views.
Governance belongs on the data asset itself, not on each downstream copy.

\section{Managed Spark and Hadoop Modernization}
\index{Dataproc!equivalence}\index{Spark!managed}%
Chapter~9's modernization story---decouple storage from compute, make clusters ephemeral, let a
serverless layer absorb administration---converges across clouds.

\begin{center}
\footnotesize
\begin{tabular}{@{}p{2.8cm}p{3.0cm}p{3.0cm}p{3.4cm}@{}}
\toprule
\textbf{Capability} & \textbf{GCP Dataproc} & \textbf{AWS EMR} & \textbf{Databricks / Azure} \\
\midrule
Managed Spark/Hadoop & Dataproc clusters & EMR on EC2 & Databricks Runtime (Photon) / HDInsight, Synapse \\
Serverless Spark & Dataproc Serverless & EMR Serverless & Serverless compute / Synapse serverless \\
Object storage for HDFS & Cloud Storage connector & S3 (EMRFS) & Unity Catalog volumes / ADLS (ABFS) \\
Cheap workers & Preemptible secondaries & Spot task nodes & Spot instances / low-priority nodes \\
Vectorized engine & (vanilla Spark) & EMR runtime (partial) & Photon / Fabric native engine \\
Cluster lifecycle & Ephemeral per job & Transient clusters & Job clusters / job-scoped pools \\
\bottomrule
\end{tabular}
\end{center}

\noindent \textbf{Snowflake}'s Snowpark absorbs DataFrame workloads without a cluster;
\textbf{MongoDB}'s aggregation framework answers smaller questions cluster-free. The migration
question is always: which workloads need the full Spark ecosystem, and which rewrite into SQL or
serverless engines entirely.

\section{Visual, Low-Code ETL}
\index{Data Fusion!equivalence}\index{low-code ETL!equivalence}%
Chapter~10's topic---visual ETL for analysts and integration engineers---targets the same persona
on every cloud.

\begin{center}
\footnotesize
\begin{tabular}{@{}p{2.8cm}p{3.0cm}p{3.0cm}p{3.4cm}@{}}
\toprule
\textbf{Capability} & \textbf{GCP Data Fusion} & \textbf{AWS} & \textbf{Azure / Fabric} \\
\midrule
Visual ETL studio & Data Fusion Studio & Glue Studio & Data Factory / Fabric Data Factory \\
Interactive cleaning & Wrangler & Glue DataBrew & Power Query / Dataflow Gen2 \\
Execution engine & Spark on ephemeral Dataproc & Glue Spark runtime & Spark / Data Factory IR \\
Extensibility & Plugin Hub (CDAP) & Glue connectors & Linked services \& connectors \\
Scheduling & Built-in schedules + triggers & Glue triggers / EventBridge & ADF triggers / Fabric schedules \\
\bottomrule
\end{tabular}
\end{center}

\noindent \textbf{Databricks} covers this niche with Lakeflow Declarative Pipelines and
partner-built connectors; \textbf{Snowflake} users reach for Matillion, Fivetran, or dbt;
\textbf{MongoDB} users lean on Atlas Data Federation and Compass. The shared trade-off: visual
tools democratize pipeline building but can hide cost and semantics behind a friendly canvas.

\section{Change Data Capture}
\index{CDC!equivalence}\index{Datastream!equivalence}%
Chapter~11's pattern is universal: read the database's transaction log, not its tables, and
deliver ordered insert/update/delete events to analytical targets.

\begin{center}
\footnotesize
\begin{tabular}{@{}p{2.8cm}p{3.0cm}p{3.0cm}p{3.4cm}@{}}
\toprule
\textbf{Capability} & \textbf{GCP Datastream} & \textbf{AWS DMS} & \textbf{Databricks / Azure} \\
\midrule
CDC engine & Log-based, serverless & DMS CDC tasks & Lakeflow Connect / ADF CDC, SQL CDC \\
PostgreSQL source & Logical replication (pgoutput) & wal2json / test\_decoding & Partner tools / PG logical replication \\
Oracle source & LogMiner / binary reader & LogMiner / binary reader & Partner or self-managed \\
Managed targets & Cloud Storage, BigQuery & S3, Redshift, Kafka & Delta tables / ADLS, Synapse \\
CDC apply & BigQuery MERGE (staging) & Glue/EMR MERGE & \texttt{AUTO CDC} / T-SQL MERGE \\
Backfill & Built-in initial snapshot & Full load + CDC & Snapshot + CDC \\
\bottomrule
\end{tabular}
\end{center}

\noindent Open-source \textbf{Debezium} on Kafka is the self-managed counterpart everywhere;
\textbf{Snowflake} consumes CDC through Snowpipe and streams; \textbf{MongoDB} has change streams
built into the database. The architectural constant: the transaction log is the source of truth,
and every destination is an eventually-consistent projection of it.

\section{Workflow Orchestration}
\index{orchestration!equivalence}\index{Composer!equivalence}%
Chapter~12's managed Airflow is the most portable orchestration story in cloud data
engineering---the same DAG Python runs on every platform.

\begin{center}
\footnotesize
\begin{tabular}{@{}p{2.8cm}p{3.0cm}p{3.0cm}p{3.4cm}@{}}
\toprule
\textbf{Capability} & \textbf{GCP Composer} & \textbf{AWS MWAA} & \textbf{Databricks / Azure} \\
\midrule
Orchestrator & Composer 2 (GKE) & MWAA (ECS-based) & Workflows (jobs + tasks) / ADF managed Airflow \\
Parameters & Airflow variables, XComs & SFN input/output & Job params + task values / pipeline params \\
Conditional flow & Branch operators & Choice / Map states & if/else, for-each tasks / If Condition \\
Partial recovery & Retries, clears & Redrive / retries & Repair run / rerun from activity \\
Secret storage & Secret Manager & Secrets Manager & Secret scopes / Key Vault \\
DAG storage & Cloud Storage bucket & S3 bucket synced & Workspace files, Asset Bundles / Git \\
\bottomrule
\end{tabular}
\end{center}

\noindent \textbf{Snowflake Tasks} and \textbf{dbt Cloud} schedules keep orchestration in-warehouse
for SQL-only pipelines; \textbf{Fabric pipelines} target low-code orchestration; \textbf{Dagster}
and \textbf{Prefect} offer asset-centric alternatives. Airflow's task-centric DAG model remains
the lingua franca the others are compared against.

\section{Messaging and Event Ingestion}
\index{Pub/Sub!equivalence}\index{messaging!equivalence}%
Chapter~13's managed messaging shares one mental model---producers publish to named channels,
consumers subscribe, the broker owns durability---with different knobs for ordering, replay, and
cost.

\begin{center}
\footnotesize
\begin{tabular}{@{}p{2.8cm}p{3.0cm}p{3.0cm}p{3.4cm}@{}}
\toprule
\textbf{Capability} & \textbf{GCP Pub/Sub} & \textbf{AWS} & \textbf{Azure / Kafka} \\
\midrule
Managed pub/sub & Pub/Sub (global) & SNS + SQS / EventBridge & Service Bus / Event Grid \\
Kafka-compatible & Managed Service for Kafka & MSK & Event Hubs Kafka surface / MSK, Confluent \\
Dead-lettering & Dead-letter topics & SQS DLQs & Service Bus dead-letter queues \\
Replay & Snapshots + seek & EventBridge archive+replay & Event Hubs offsets / Kafka retention \\
Schema enforcement & Avro / Proto topic schemas & EventBridge schemas & Schema Registry / Azure Schema Registry \\
Cost-profiled tier & Pub/Sub Lite (zonal) & Kinesis on-demand vs provisioned & Event Hubs throughput units / partitions \\
\bottomrule
\end{tabular}
\end{center}

\noindent \textbf{Databricks} and \textbf{Snowflake} sit downstream of these buses (Spark Kafka
source; Snowpipe Streaming, Kafka Connector); \textbf{MongoDB} consumes via Atlas Stream
Processing. Pub/Sub differs deliberately from Kafka: global by default, partition-free for the
user, priced per GiB delivered rather than per provisioned broker.

\section{Stream Processing Engines and Semantics}
\index{Dataflow!equivalence}\index{stream processing!equivalence}\index{watermarks!equivalence}%
Chapters~14 and~15 teach Beam's unified model; every serious engine answers the Dataflow Model's
four questions with different defaults and vocabulary.

\begin{center}
\footnotesize
\begin{tabular}{@{}p{2.8cm}p{2.9cm}p{2.9cm}p{3.6cm}@{}}
\toprule
\textbf{Concept} & \textbf{Beam / Dataflow} & \textbf{Apache Flink} & \textbf{Spark Structured Streaming (Databricks)} \\
\midrule
Event-time windows & Fixed / sliding / session & Tumbling / hopping / session & \texttt{window()}, session windows \\
Completeness & Watermark & Watermark & \texttt{withWatermark} \\
Late data & Allowed lateness + triggers & Allowed lateness + side outputs & Watermark delay \\
Per-key state & State \& Timers API & Keyed state + timers & \texttt{flatMapGroupsWithState} \\
Dedup & Stateful DoFn & Keyed state + TTL & \texttt{dropDuplicates} + watermark \\
Exactly-once & Shuffle + Storage Write API & Checkpointing + 2PC sinks & Checkpoint + idempotent sink \\
Batch unification & One model, both modes & Flink batch mode & \texttt{availableNow} trigger \\
\bottomrule
\end{tabular}
\end{center}

\noindent Managed runners exist on every cloud: Dataflow (GCP), Managed Service for Apache Flink
and Kinesis Data Analytics (AWS), Stream Analytics and Fabric eventstreams (Azure).
\textbf{Snowflake} lands micro-batches with simpler semantics; \textbf{MongoDB} change streams feed
processors that implement these semantics themselves. Beam's vocabulary---watermark, trigger,
allowed lateness---is the common reference even when the runtime differs \cite{akidau2015dataflow}.

\section{Governed BI and Semantic Layers}
\index{Looker!equivalence}\index{semantic layer!equivalence}%
Chapter~16's governance question---are metrics defined once in a semantic layer, or independently
in every dashboard?---confronts every platform.

\begin{center}
\footnotesize
\begin{tabular}{@{}p{2.8cm}p{3.0cm}p{3.0cm}p{3.4cm}@{}}
\toprule
\textbf{Capability} & \textbf{GCP Looker} & \textbf{AWS} & \textbf{Azure / Databricks} \\
\midrule
Governed BI + semantics & Looker (LookML) & QuickSight datasets & Power BI semantic models / AI/BI + metrics \\
Self-serve dashboards & Looker Studio & QuickSight & Power BI / AI/BI dashboards \\
Model as code & LookML in Git & QuickSight APIs & TMDL / Genie instructions \\
Natural-language BI & Gemini in Looker & Amazon Q in QuickSight & Copilot / Genie spaces \\
In-memory acceleration & BI Engine & SPICE & Import / DirectLake / SQL warehouse cache \\
Reverse ETL & Scheduled queries + partners & Census, Hightouch & Partner tools \\
\bottomrule
\end{tabular}
\end{center}

\noindent \textbf{Snowflake} pairs with the same BI tools over its warehouses; \textbf{MongoDB
Charts} serves operational analytics on Atlas. The portable insight: a semantic layer is a
\emph{code artifact}---versioned, reviewed, tested---and organizations that treat it as such are
the ones where ``Gross Revenue'' means one thing.

\section{In-Warehouse Machine Learning and Feature Stores}
\index{BigQuery ML!equivalence}\index{feature store!equivalence}%
Chapter~17's in-warehouse ML eliminates the extract-train-reload loop; every major warehouse ships
a variant, and every serious ML platform implements a feature store as the contract between
training and serving.

\begin{center}
\footnotesize
\begin{tabular}{@{}p{2.8cm}p{3.0cm}p{3.0cm}p{3.4cm}@{}}
\toprule
\textbf{Capability} & \textbf{GCP} & \textbf{Snowflake} & \textbf{Databricks / AWS} \\
\midrule
SQL-native training & BigQuery ML \texttt{CREATE MODEL} & Snowpark ML / ML functions & SQL + MLlib / Redshift ML \\
Forecasting & ARIMA\_PLUS & Forecasting ML function & AutoML / SageMaker forecasting \\
Feature registry & Vertex AI Feature Store & Snowflake Feature Store & Feature Engineering in UC / SageMaker FS \\
Offline store & BigQuery tables & Snowflake tables & Delta tables (UC) / S3 + Glue catalog \\
Online serving & Feature Store online & (hybrid tables) & Feature serving endpoints / ElastiCache \\
Point-in-time join & Training-set builders & ASOF joins & \texttt{create\_training\_set} / event-time joins \\
\bottomrule
\end{tabular}
\end{center}

\noindent \textbf{Azure / Fabric} fields Synapse and Fabric ML with Spark MLlib; \textbf{MongoDB
Atlas} exports to analytical stores for ML. The shared virtue: governance, lineage, and access
control of the data platform extend automatically to training data, features, and models.

\section{ML Platforms, Pipelines, and Serving}
\index{Vertex AI!equivalence}\index{MLOps!equivalence}\index{model registry!equivalence}%
Chapters~18 and~19's managed ML platforms converge on the same lifecycle---train, register,
deploy, explain, monitor, retrain---behind different consoles.

\begin{center}
\footnotesize
\begin{tabular}{@{}p{2.8cm}p{3.0cm}p{3.0cm}p{3.4cm}@{}}
\toprule
\textbf{Capability} & \textbf{GCP Vertex AI} & \textbf{AWS SageMaker} & \textbf{Databricks / Azure ML} \\
\midrule
AutoML (tabular) & AutoML Tabular & Autopilot & AutoML / Automated ML \\
Experiment tracking & Vertex Experiments & SageMaker Experiments & Managed MLflow / Azure ML runs \\
Model registry & Model Registry & Model Registry & Unity Catalog registry / MLflow registry \\
Pipelines & Vertex Pipelines (KFP/TFX) & SageMaker Pipelines & Workflows + MLflow / Azure ML pipelines \\
Online serving & Endpoints & Endpoints & Model Serving (serverless) / online endpoints \\
Batch scoring & Batch prediction & Batch transform & Spark UDF from registry / batch endpoints \\
Monitoring & Model Monitoring & Model Monitor & Lakehouse Monitoring / Azure ML monitoring \\
Request logging & Endpoint logs & Data capture & Inference tables (Delta) / App Insights \\
\bottomrule
\end{tabular}
\end{center}

\noindent \textbf{Snowflake} embeds Snowpark ML and Cortex for in-warehouse inference;
\textbf{MongoDB} users serve models adjacent to Atlas. For the data engineer, the constant is the
interface contract: features in, predictions out, lineage and governance around both---and
retraining as an \emph{event-driven reaction to evidence}, not a calendar entry.

\section{Pretrained AI APIs and Generative AI}
\index{AI APIs!equivalence}\index{generative AI!equivalence}\index{RAG!equivalence}%
Chapter~20's pretrained APIs and the retrieval-augmented generation (RAG) pattern are the fastest
paths from unstructured data to structured signal on every cloud.

\begin{center}
\footnotesize
\begin{tabular}{@{}p{2.8cm}p{3.0cm}p{3.0cm}p{3.4cm}@{}}
\toprule
\textbf{Capability} & \textbf{GCP} & \textbf{AWS} & \textbf{Azure / Databricks} \\
\midrule
Image analysis / OCR & Vision API & Rekognition / Textract & AI Vision / Document Intelligence \\
Text NLP & Natural Language API & Comprehend & Azure AI Language \\
Translation & Translation API & Translate & Azure Translator \\
Speech & Speech-to-Text & Transcribe & Azure Speech \\
Foundation models & Gemini (Vertex AI) & Bedrock & Azure OpenAI / Foundation Model APIs \\
Vector index & Vertex Vector Search & OpenSearch / Kendra & AI Search / Mosaic AI Vector Search \\
RAG orchestration & Vertex RAG Engine & Bedrock Knowledge Bases & Azure AI RAG / LangChain + MLflow \\
\bottomrule
\end{tabular}
\end{center}

\noindent \textbf{MongoDB Atlas Vector Search} and \textbf{Snowflake Cortex} bring embeddings and
LLM functions into the data layer. The engineering constants: per-call pricing demands cost
forecasting, and external calls inside distributed pipelines demand batching and backoff.

\section{Security: Perimeters, Keys, and Identity}
\index{VPC Service Controls!equivalence}\index{CMEK!equivalence}%
Chapter~21's perimeter-based exfiltration defense and key management have direct counterparts on
every cloud.

\begin{center}
\footnotesize
\begin{tabular}{@{}p{2.8cm}p{3.0cm}p{3.0cm}p{3.4cm}@{}}
\toprule
\textbf{Capability} & \textbf{GCP} & \textbf{AWS} & \textbf{Azure / Databricks} \\
\midrule
Service perimeter & VPC Service Controls & SCPs + endpoint policies & Private Link + policy / IP access lists \\
Private API access & Private Google Access / PSC & VPC endpoints (PrivateLink) & Private Link / private connectivity \\
Key management & Cloud KMS (CMEK/CSEK) & KMS (CMK/BYOK) & Key Vault (CMK/BYOK) / CMK \\
Identity federation & Workload Identity Federation & IAM Roles Anywhere / IRSA & Workload identity federation \\
Machine identity & Service accounts & IAM roles & Managed identity / service principal + OAuth \\
Audit trail & Cloud Audit Logs & CloudTrail & Activity Log / \texttt{system.access.audit} \\
\bottomrule
\end{tabular}
\end{center}

\noindent \textbf{Snowflake} implements perimeter control with network policies and Tri-Secret
Secure; \textbf{MongoDB Atlas} with private endpoints and field-level encryption. The shared
mental model: IAM answers \emph{who may}, perimeters answer \emph{from where and to where}, and
keys answer \emph{who can read it at all}.

\section{Governance Fabrics, Lineage, and Data Quality}
\index{Dataplex!equivalence}\index{governance!equivalence}\index{data quality!equivalence}%
Chapter~22's governance fabric---metadata, quality, and policy attached to data wherever it
lives---differs per platform in scope and tightness of integration.

\begin{center}
\footnotesize
\begin{tabular}{@{}p{2.8cm}p{3.0cm}p{3.0cm}p{3.4cm}@{}}
\toprule
\textbf{Capability} & \textbf{GCP Dataplex} & \textbf{AWS} & \textbf{Azure / Databricks} \\
\midrule
Governance fabric & Dataplex (lakes/zones) & Lake Formation & Purview / Unity Catalog \\
Catalog \& search & Data Catalog (in Dataplex) & Glue Data Catalog & Purview / UC metastore \\
Lineage & Dataplex lineage & Glue (partial) / SageMaker & Purview lineage / UC auto-lineage \\
Auto data quality & AutoDQ scans & Glue Data Quality (Deequ) & Purview / Lakehouse Monitoring, DLT expectations \\
Data mesh units & Lakes per domain & Accounts per domain & Fabric domains / catalogs per domain \\
\bottomrule
\end{tabular}
\end{center}

\noindent \textbf{Snowflake} governs through Horizon (tags, masking, lineage) inside its boundary;
\textbf{MongoDB} relies on schema validation and Atlas roles. The mesh principle is
platform-agnostic: domains own data products; the platform provides federated
guardrails---discovery, quality, and policy enforcement---without centralizing every decision.

\section{Sensitive-Data Discovery and De-identification}
\index{DLP!equivalence}\index{de-identification!equivalence}%
Chapter~23's sensitive-data discovery and de-identification exist on every platform, though GCP's
Sensitive Data Protection is unusually deep as a standalone service.

\begin{center}
\footnotesize
\begin{tabular}{@{}p{2.8cm}p{3.0cm}p{3.0cm}p{3.4cm}@{}}
\toprule
\textbf{Capability} & \textbf{GCP} & \textbf{AWS} & \textbf{Azure / Databricks} \\
\midrule
PII discovery & Sensitive Data Protection & Macie & Purview Information Protection / UC discovery \\
De-identification & DLP transforms (mask, tokenize, FPE) & Partner / Lambda transforms & Purview / ADF masking / UC masks \\
Warehouse masking & Policy tags + routines & Lake Formation + UDFs & Dynamic data masking / column masks \\
Format-preserving & DLP FFX/FPE & (partner) & (partner) \\
Risk analysis & k-anonymity / l-diversity & (partner) & (partner) \\
\bottomrule
\end{tabular}
\end{center}

\noindent \textbf{Snowflake} provides masking policies and tag-based protection; \textbf{MongoDB}
provides field-level encryption. The portable compliance lesson: deletion is not
erasure---time travel, backups, and replicas retain copies, and only cryptographic erasure or
exhaustive replay satisfies a strict right-to-erasure obligation.

\section{Infrastructure as Code, SQL Pipelines, and CI/CD}
\index{Terraform!equivalence}\index{CI/CD!equivalence}\index{DataOps!equivalence}%
Chapter~24's disciplines---infrastructure and transformations as code---are platform-portable;
only the resource names change.

\begin{center}
\footnotesize
\begin{tabular}{@{}p{2.8cm}p{3.0cm}p{3.0cm}p{3.4cm}@{}}
\toprule
\textbf{Capability} & \textbf{GCP} & \textbf{AWS} & \textbf{Azure / Databricks} \\
\midrule
IaC & Terraform google provider / Deployment Manager & Terraform aws / CloudFormation & Terraform azurerm / Bicep / Asset Bundles \\
SQL pipelines & Dataform (SQLX, assertions) & dbt on Redshift / Athena & dbt on Synapse / Databricks SQL, DLT \\
CI/CD & Cloud Build + Workload Identity & CodePipeline / GitHub Actions & Azure DevOps / GitHub Actions \\
State backend & GCS backend with locking & S3 + DynamoDB locking & Azure Storage / Terraform remote state \\
Validate before deploy & \texttt{terraform plan} & cfn-lint / CDK synth & what-if / \texttt{bundle validate} \\
Cost telemetry & Billing export to BigQuery & CUR 2.0 & Cost Management / \texttt{system.billing.usage} \\
Exam analog & Professional Data Engineer & Data Engineer Associate & DP-203/DP-600 / Databricks Data Engineer \\
\bottomrule
\end{tabular}
\end{center}

\noindent \textbf{Snowflake} practitioners reach for dbt and Terraform's Snowflake provider;
\textbf{MongoDB} for the Atlas provider and schema-versioned migrations. The DataOps constant:
the merge is the deploy, and the pipeline is the only writer to production.

\section{Databricks Equivalents at a Glance}
\index{Databricks!equivalence summary}%
Because Databricks runs on all three hyperscaler clouds, it deserves a single consolidated mapping
against this book's GCP services, harvested from the companion Databricks book's cross-cloud
boxes.

\begin{center}
\footnotesize
\begin{tabular}{@{}p{4.6cm}p{4.6cm}p{3.6cm}@{}}
\toprule
\textbf{GCP service / concept} & \textbf{Databricks equivalent} & \textbf{Book chapter} \\
\midrule
GCP project + console & Workspace + account console & Ch.~0 / DB Ch.~1 \\
Cloud Storage (lake layout) & Volumes + external locations (Delta on GCS/S3/ADLS) & Ch.~1 / DB Ch.~2 \\
BigQuery table & Delta table in Unity Catalog & Ch.~5 / DB Ch.~2 \\
Partitioning + clustering & Partitioning, \texttt{ZORDER BY}, liquid clustering & Ch.~6 / DB Ch.~5 \\
BigQuery reservations / slots & SQL warehouses (serverless / pro / classic) & Ch.~6 / DB Ch.~6 \\
BigLake governance & Unity Catalog (grants, row filters, masks) & Ch.~8 / DB Ch.~8 \\
Looker / Looker Studio & AI/BI dashboards + Genie & Ch.~16 / DB Ch.~7 \\
Dataproc (managed Spark) & Databricks Runtime clusters (Photon) & Ch.~9 / DB Ch.~4 \\
Data Fusion (visual ETL) & Lakeflow Declarative Pipelines (DLT) & Ch.~10 / DB Ch.~10 \\
Datastream (CDC) & Lakeflow Connect + \texttt{AUTO CDC} & Ch.~11 / DB Ch.~11 \\
Composer (Airflow) & Workflows (jobs, tasks, repair runs) & Ch.~12 / DB Ch.~12 \\
Pub/Sub + Dataflow (streaming) & Kafka sources + Structured Streaming & Ch.~13--15 / DB Ch.~13--16 \\
Storage Write API upserts & \texttt{foreachBatch} + \texttt{MERGE} & Ch.~15 / DB Ch.~16 \\
Vertex Experiments / Registry & Managed MLflow + UC model registry & Ch.~18 / DB Ch.~17 \\
Vertex Feature Store & Feature Engineering in UC + serving & Ch.~17, 19 / DB Ch.~18 \\
Vertex endpoints / monitoring & Model Serving + Lakehouse Monitoring & Ch.~18, 20 / DB Ch.~20 \\
Vertex Vector Search / RAG & Mosaic AI Vector Search + Foundation Models & Ch.~20 / DB Ch.~19 \\
VPC SC, KMS, Workload Identity & Private connectivity, CMK, service principals & Ch.~21 / DB Ch.~21 \\
Dataplex (lineage, audit) & UC lineage + \texttt{system.access.audit} & Ch.~22 / DB Ch.~22 \\
Analytics Hub (sharing) & Delta Sharing (open protocol) & Ch.~8 / DB Ch.~22 \\
Terraform + Cloud Build & Asset Bundles + Terraform + CI/CD & Ch.~24 / DB Ch.~23 \\
\bottomrule
\end{tabular}
\end{center}

\noindent The recurring contrast: GCP splits the platform into purpose-built managed services with
separate consoles, while Databricks consolidates the same capabilities around one runtime, one
table format (Delta), and one governance plane (Unity Catalog). Neither is simpler---they
distribute complexity differently.

\begin{tipbox}
Use this appendix as a study aid for the book's lockstep, cross-cloud approach: when you learn a
concept on one platform, immediately locate its equivalent here and predict how the other five
would express the same idea. That habit turns six separate vocabularies into one mental model.
\end{tipbox}
